% Packages
\usepackage{color,graphicx}
\usepackage[spanish,activeacute]{babel}
\usepackage{amsmath,amsthm,amssymb}
\usepackage{pgf,tikz}
\usetikzlibrary{arrows,positioning,calc,intersections,petri}
\usetikzlibrary{lindenmayersystems,decorations}
\usepackage{hyperref}
\usepackage{multicol}
\usepackage{bm}
\usepackage{centernot}
\usepackage{nicematrix,tikz}
\usepackage{epsfig,mathrsfs}
\usepackage{amsfonts,xspace,colortbl,enumerate}
\usepackage[utf8]{inputenc}

% itemize
\usepackage{enumitem}
\setlist{nolistsep}

% Dimensions
\setlength{\oddsidemargin}{0cm}
\setlength{\evensidemargin}{0cm} 
\setlength{\textwidth}{17cm}
\setlength{\textheight}{23cm} 
\setlength{\topmargin}{-1cm}

\addtolength{\voffset}{0cm}

\parindent= 0cm
\parskip=1em

\font\bff=cmbx10 at 10truept
\font\lg=cmdunh10 at 10truept
\font\bl=cmss10 at 10truept

\linespread{1.1}
% ----------------------------------------------------------------
\vfuzz2pt % Don't report over-full v-boxes if over-edge is small
\hfuzz2pt % Don't report over-full h-boxes if over-edge is small

% useful for formulas
% Conjuntos de numeros
\newcommand{\R}{\mathbb{R}}
\newcommand{\N}{\mathbb{N}}
\newcommand{\Z}{\mathbb{Z}}
\newcommand{\C}{\mathbb{C}}
\newcommand{\Q}{\mathbb{Q}}
\newcommand{\K}{\mathbb{K}}
% En demostraciones
\newcommand{\gdw}{\Leftrightarrow}
\newcommand{\dw}{\Rightarrow}
% Functions
\newcommand{\To}{\longrightarrow}
% derivative
\newcommand{\de}{\mbox{d}}
% distance
\newcommand{\dist}{\mbox{d}}
% linear operators
\newcommand{\linop}{{\mathcal L}}
% image space of a matrix
\newcommand{\myim}{\mbox{im}\,}
\newcommand{\ssi}{si y solo si }
\newcommand\sbullet[1][.5]{\mathbin{\vcenter{\hbox{\scalebox{#1}{$\bullet$}}}}}

% theorem, lemma, example, corollary, proposition, ...
\theoremstyle{definition}
\newtheorem{thm}{Teorema}[chapter]
\theoremstyle{definition}
\newtheorem{cor}{Corolario}[thm]
\theoremstyle{definition}
\newtheorem{lem}[thm]{Lema}
\theoremstyle{definition}
\newtheorem{defn}[thm]{Definici\'on}
\newtheorem{prop}[thm]{Proposici\'on}


\theoremstyle{plain}
\newtheorem{ejemplo}[thm]{Ejemplo}
\theoremstyle{plain}
\newtheorem{ejercicio}[thm]{Ejercicio}
\theoremstyle{plain}
\newtheorem{rem}[thm]{Observaci\'on}

\newenvironment{prueba}{\begin{proof}[\textbf{Demostraci\'on}]}{\end{proof}}

\numberwithin{equation}{chapter}

\definecolor{reddish}{rgb}{.8,.15,.2}

%--------------------------------------------
% Recuadros para lemas, teoremas,
% observaciones y definiciones
\usepackage{caption}
\usepackage{parskip}
\usepackage{tcolorbox}
\tcbuselibrary{theorems}


% Colores A. Hopper
% verde: blue!35!black (definiciones)
% rojo: red!90!black (observaciones)
% azul: (ejemplos)

\definecolor{dgreen}{rgb}{0.0,0.6,0.0}
\definecolor{turco}{rgb}{0,.5,.5}
\definecolor{reddish}{rgb}{.8,.1,.1}
\definecolor{magic}{cmyk}{0,.5,0,.5}
\definecolor{blueish}{rgb}{.1,.1,.8}
\definecolor{greenish}{rgb}{.1,.8,.1}
\definecolor{antiquefuchsia}{rgb}{0.57, 0.36, 0.51}%
\definecolor{amaranth}{rgb}{0.9, 0.17, 0.31}
\definecolor{awesome}{rgb}{1.0, 0.13, 0.32}
\definecolor{cadmiumgreen}{rgb}{0.0, 0.42, 0.24}
\definecolor{brightmaroon}{rgb}{0.76, 0.13, 0.28}
\definecolor{cadmiumred}{rgb}{0.89, 0.0, 0.13}
\definecolor{carminered}{rgb}{1.0, 0.0, 0.22}
\definecolor{cobalt}{rgb}{0.0, 0.28, 0.67}

\newcommand{\az}[1]{\textcolor{blueish}{#1}}
\newcommand{\ro}[1]{\textcolor{reddish}{#1}}
\newcommand{\tu}[1]{\textcolor{turco}{#1}}
\newcommand{\na}[1]{\textcolor{brightmaroon}{#1}}
\newcommand{\aw}[1]{\textcolor{cobalt}{#1}}

\newcommand*\circled[1]{\tikz[baseline=(char.base)]{
		\node[shape=circle,draw,inner sep=2pt] (char) {#1};}}

%\newtcbtheorem[number within=chapter]{defn}{Definici\'on}%
%{fonttitle=\bfseries}{it}

%\newtcbtheorem[use counter from = defn]{lem}{Lema}%
%{fonttitle=\bfseries}{it}

%\newtcbtheorem[use counter from = defn]{rem}{Observaci\'on}%
%{fonttitle=\bfseries}{it}
%
%\newtcbtheorem[use counter from = defn]{ejemplo}{Ejemplo}%
%{fonttitle=\bfseries}{it}
%
%\newtcbtheorem[use counter from = defn]{cor}{Corolario}%
%{fonttitle=\bfseries}{it}
%
%\newtcbtheorem[use counter from = defn]{thm}{Teorema}%
%{fonttitle=\bfseries}{it}



%--------------------------------------------
\hypersetup{ pdfborder = {0 0 0} }

\hypersetup{pdfstartview={XYZ null null 1.17}}

\voffset 0 cm

%----------------------------------------------
% Alternativa para tener listas en columnas con filas alineadas
\usepackage{tasks}

\usepackage{circuitikz}

%\pagestyle{empty}
\parindent=0cm
\parskip=1em

% Conjuntos de numeros
\newcommand{\calL}{\mathcal{L}}
\newcommand{\sen}{\mathup{sen}}
% En demostraciones
\newcommand{\un}{{\mathcal U}}
% cardinal de un conjunto
\def\card{\textup{card}}

\newenvironment{solucion}[1]
{\begin{proof}[{\bf Soluci\'on a ejercicio #1}]}{\end{proof}}
% colores
\definecolor{ffqqqq}{rgb}{1,0,0}
\definecolor{qqqqff}{rgb}{0,0,1}
\definecolor{uququq}{rgb}{0.25,0.25,0.25}
\definecolor{qqffqq}{rgb}{0,1,0}
\definecolor{uuuuuu}{rgb}
{0.26666666666666666,0.26666666666666666,0.26666666666666666}

\newtheorem*{pp}{Principio de Inducción}


\pgfdeclaredecoration{quasi-sirpinski}{do}{%
    \state{do}[width=\pgfdecoratedinputsegmentlength, next state=do]{%
        \pgfpathmoveto{\pgfpointpolar{-60}{\pgfdecoratedinputsegmentlength/2}}%
        \pgfpathlineto{\pgfpointorigin}%
        \pgfpathlineto{\pgfpoint{\pgfdecoratedinputsegmentlength/2}{0pt}}%
        \pgfpathclose%
    }
}
